\chapter{Введение}
\label{ch:intro}

% исторически введение пишется самым последним, после определения цели и расчётов

% про изменение параметров, получилась неэффективная установка, взяли более широкую трубку

В производстве электронных компонентов, таких как полупроводниковые микросхемы, важную роль играют вакуумные системы. Например, при нанесении тонкоплёночных покрытий методом магнетронного напыления требуется быстрое достижение высокого вакуума перед началом процесса. Если откачка происходит слишком медленно, это приводит к увеличению времени выполнения технологического процесса, окислению подложек остаточными газами и, как следствие, к браку продукции. Таким образом, контроль скорости откачки вакуума напрямую влияет на качество продукции и экономическую эффективность производства.

Первым методом контроля скорости откачки является эмпирический подбор параметров насоса и времени откачки. Например, оператор может включать вакуумный насос на фиксированный промежуток времени, основываясь на предыдущем опыте, без точного контроля скорости откачки. Однако такой подход не учитывает возможные изменения в системе, такие как утечки, загрязнения и износ насоса, что приводит к нестабильности остаточного давления и перерасходу ресурсов.

Более эффективным решением является использование  методов, позволяющих измерять скорость откачки в реальном времени. Такой подход обеспечивает точное управление процессом, своевременное обнаружение неисправностей и адаптацию к изменяющимся условиям. Целью работы было определение скорости откачки по ухудшению и улучшению вакуума с использованием вакууметра.